\documentclass[../../../include/open-logic-section]{subfiles}

\begin{document}

\olfileid[fr]{sth}{ord-arithmetic}{mult}
\olsection{Multiplication ordinale}

Passons à la multiplication ordinale, que nous aborderons
à peu près comme l'addition. Supposons donc que nous
voulions multiplier $\alpha$ par~$\beta$. On peut
imaginer une grille rectangulaire, de largeur $\alpha$
et de hauteur $\beta$ : le produit de $\alpha$ et de
$\beta$ est obtenu en parcourant chaque rangée, puis
la suivante\ldots jusqu'à avoir parcouru la grille
entière. Autrement dit, le produit de $\alpha$ et
de $\beta$ s'obtient en remplaçant \emph{chaque}
élément de $\beta$ par une copie de $\alpha$.

Pour formaliser cette idée, employons simplement
l'ordre lexicographique inverse sur le produit
cartésien de $\alpha$ et de $\beta$ :

\begin{defn}
Pour tous ordinaux $\alpha, \beta$, leur produit est
$\alpha \ordtimes \beta =
\ordtype{\alpha \times \beta, \rlexless}$.
\end{defn}
\noindent
Il nous faut à nouveau vérifier que cette définition
est correcte :

\begin{lem}\ollabel{ordtimeslessiswo}
Pour tous ordinaux $\alpha$ et $\beta$,
$\tuple{\alpha \times \beta, \rlexless}$
est un bon ordre.
\end{lem}

\begin{proof}
Exactement comme pour le \olref[add]{ordsumlessiswo}.
\end{proof}
\noindent
Il n'est pas difficile non plus de démontrer que la
multiplication satisfait les équations suivantes :

\begin{lem}\ollabel{ordtimesrecursion}
Pour tous ordinaux $\alpha, \beta$ :
\begin{align*}
	\alpha \ordtimes 0 &= 0\\
	\alpha \ordtimes (\beta \ordplus 1)
	&= (\alpha \ordtimes \beta) \ordplus \alpha\\
	\alpha \ordtimes \beta
	&= \bigcup_{\delta < \beta}
	(\alpha \ordtimes \delta)
	&& \text{si $\beta$ est un ordinal limite}.\footnotemark
\end{align*}
\footnotetext{Note de l'édition française : la borne supérieure stricte de l'original donnerait $1$ pour $0\ordtimes\beta$ lorsque $\beta$ est limite. L'union des produits antérieurs vaut $0$, conformément à la définition par produit cartésien, et convient aussi aux facteurs non nuls.}\end{lem}

\begin{proof}
La démonstration est laissée en exercice.
\end{proof}

Comme pour l'addition, nous aurions d'ailleurs pu
définir la multiplication ordinale à l'aide de ces
équations de récursion, sans en donner directement
une définition. Comme pour l'addition encore,
certaines propriétés nous sont familières :

\begin{lem}\ollabel{ordinalmultiplicationisnice}
Si $\alpha, \beta, \gamma$ sont des ordinaux, alors :
\begin{enumerate}
	\item\ollabel{ordtimes1}
	si $\alpha \neq 0$ et $\beta < \gamma$,
	alors $\alpha \ordtimes \beta <
	\alpha \ordtimes \gamma$ ;
	\item\ollabel{ordtimes2}
	si $\alpha \neq 0$ et
	$\alpha \ordtimes \beta =
	\alpha\ordtimes\gamma$, alors $\beta = \gamma$ ;
	\item\ollabel{ordtimes3}
	$\alpha \ordtimes (\beta \ordtimes \gamma)
	= (\alpha \ordtimes \beta) \ordtimes \gamma$ ;
	\item\ollabel{ordtimes4}
	si $\alpha \leq \beta$, alors
	$\alpha \ordtimes \gamma
	\leq \beta \ordtimes \gamma$ ;
	\item\ollabel{ordtimes5}
	$\alpha \ordtimes (\beta \ordplus \gamma)
	= (\alpha \ordtimes \beta)
	\ordplus (\alpha\ordtimes \gamma)$.
\end{enumerate}
\end{lem}

\begin{proof}
La démonstration est laissée en exercice.
\end{proof}

Vous pouvez démontrer d'autres résultats, ou les
chercher dans la littérature, à votre guise.
Mais, compte tenu de la
\olref[ord-arithmetic][add]{ordsumnotcommute},
le résultat suivant ne devrait pas surprendre :

\begin{prop}
La multiplication ordinale n'est pas commutative :
$2 \ordtimes \omega = \omega <
\omega \ordtimes 2$.
\end{prop}

\begin{proof}
$2 \ordtimes \omega =
\supstrict_{n < \omega}(2\ordtimes n)
= \omega \in
\supstrict_{n < \omega}(\omega \ordplus n)
= \omega \ordplus \omega
= \omega \ordtimes 2$.
\end{proof}
\noindent
L'explication intuitive est, là encore, assez simple.
Pour calculer $2 \ordtimes \omega$, on remplace
chaque nombre naturel par deux objets. On obtiendrait
le même type d'ordre en insérant dans les naturels
tous les nombres «~demi-entiers~», autrement dit
en considérant l'ordre usuel sur
$\Setabs{\nicefrac{n}{2}}{n \in \omega}$.
Sous cette forme, il est clair que le type d'ordre
est celui de $\omega$ lui-même. Pour calculer
$\omega \ordtimes 2$, en revanche, on place deux
copies de $\omega$ l'une après l'autre.

\begin{prob}
Démontrez le
\olref[sth][ord-arithmetic][mult]{ordtimeslessiswo},
le \olref[sth][ord-arithmetic][mult]{ordtimesrecursion}
et le
\olref[sth][ord-arithmetic][mult]{ordinalmultiplicationisnice}.
\end{prob}

\end{document}
